\documentclass{article}
\usepackage[T1]{fontenc}
\usepackage{lmodern}
\usepackage{graphicx}
\usepackage{amsmath,amssymb}
\usepackage[a4paper,margin=2cm]{geometry}
\usepackage[absolute,overlay]{textpos}
\setlength{\TPHorizModule}{1mm}
\setlength{\TPVertModule}{1mm}
\usepackage[indonesian]{babel}
\usepackage{hyperref}
\setlength{\emergencystretch}{2em}
% unit: Halaman 10

\begin{document}

% === Halaman 10 (python/native) ===

Pembangkitan Kunci Internal (key expansion)

• Kunci internal = kunci setiap putaran = subkey (round key)

• Ada 16 putaran, jadi ada 16 kunci internal: K1, K2, …, K16 

• Kunci internal dibangkitkan dari kunci eksternal (64 bit) yang diberikan oleh 

pengguna.  Di dalam dokumen resmi DES pembangkitan kunci internal disebut 

key schedule.

• Panjang kunci internal adalah 48 bit

• Gambar 3 memperlihatkan proses pembangkitan kunci internal.

10

\end{document}
